\documentclass[11pt]{article}
\usepackage[T1]{fontenc}
\usepackage[margin=1in]{geometry}
\usepackage{amsmath,amssymb,mathtools,bm}
\usepackage{booktabs,array,microtype}
\usepackage[hidelinks]{hyperref}
\usepackage{xurl}
\numberwithin{equation}{section}
\allowdisplaybreaks[2]
\setlength{\parskip}{0.35em}
\newcommand{\NS}{\mathrm{NS}}
\newcommand{\fq}{\mathfrak q}
\newcommand{\cF}{\mathcal F}
\newcommand{\cG}{\mathcal G}
\newcommand{\cC}{\mathcal C}
\newcommand{\ve}{\boldsymbol\epsilon}
\newcommand{\vh}{\boldsymbol h}
\newcommand{\vd}{\boldsymbol d}
\newcommand{\bpzbraket}[2]{\bigl\langle\!\bigl\langle #1\bigm|#2\bigr\rangle}
\DeclareMathOperator*{\Res}{Res}

\title{Sphere \texorpdfstring{$n$}{n}-Point NS Blocks in the Pillow Frame\\
and the Proposed Elliptic \texorpdfstring{$h$}{h}-Recursion}
\author{Machine note}
\date{August 30, 2026}

\begin{document}
\maketitle

\begin{abstract}
We give a self-contained recipe for bottom-component Neveu--Schwarz
sphere blocks in the comb, or open-necklace, channel. The ordinary pillow
map and its geometric conformal factor are retained. After primary
propagation powers are removed, the proposed regular part is a universal
infinite product; the singular part is determined by NS null-vector
factorization and closes the elliptic recursion. We use the conventions
of \texttt{Human Notes/SCblock.tex}, including its normalized elliptic
integral, torus coordinate, central charge, and three-point forms. The
general multipoint large-weight seed is an assumption supported by the
five- and six-point checks recorded below, not a proved cap Ward identity. This
standalone note neither inputs nor modifies the human note.
\end{abstract}

\section{Conventions and scope}
\label{sec:conventions}

We use the ordinary central charge of the human note:
\begin{align}
 [L_m,L_n]&=(m-n)L_{m+n}
       +\frac{c}{12}(m^3-m)\delta_{m+n,0},\notag\\
 [L_m,G_r]&=\left(\frac m2-r\right)G_{m+r},\qquad
 \{G_r,G_s\}=2L_{r+s}
       +\frac c3\left(r^2-\frac14\right)\delta_{r+s,0},
 \label{eq:algebra}\\
 c&=\frac32+3Q^2,\qquad Q=b+b^{-1},\qquad
 r,s\in\mathbb Z+\frac12\quad\hbox{in the algebra}.
 \notag
\end{align}
In subsequent Kac labels, $r,s$ will instead be positive integers.
We take external bottom components of intrinsically even NS primaries,
and choose even highest-weight representatives for the internal modules.
Upper external components and Ramond modules require separate seeds and
are not included here.

Put the external weights $d_1,\ldots,d_n$ at the ordered positions
\begin{equation}
 (0,z,t_1,\ldots,t_{n-4},1,\infty),\qquad
 0<z<t_1<\cdots<t_{n-4}<1,
 \label{eq:positions}
\end{equation}
initially on the real OPE sheet. Write $m=n-3$ for the number of internal
edges, with weights $h_1,\ldots,h_m$. This follows the human note's
necklace convention: $h_i$ are internal and $d_i$ external. In its
four-point subsection, the four external weights denoted there by $h_i$
are consequently relabelled $d_i$ here.

The human-note descendant and three-form conventions are
\begin{align}
 \mathbf L_{-A}\phi
 &=L_{-n_1}\cdots L_{-n_\ell}G_{-r_1}\cdots G_{-r_p}\phi,
 &(-1)^A&=(-1)^p,\notag\\
 \rho_0(\phi_1,\phi_2,\phi_3)&=1,
 &\rho_1(\phi_1,G_{-1/2}\phi_2,\phi_3)&=1,
 \label{eq:three-form-convention}\\
 (\mathbf G_h)_{AB}
 &=\bpzbraket{\mathbf L_{-A}\phi_h}{\mathbf L_{-B}\phi_h}.
 \notag
\end{align}
The relative label on $\rho_\alpha$ is total descendant parity, not the
absolute parity including the primaries. We keep the fixed-parity
ordering of the human note; the local fusion-polynomial arguments must
not be interchanged without their reflection signs.

Let $\epsilon_i\in\{0,1\}$ select levels
$\mathbb Z_{\geq0}+\epsilon_i/2$ on internal edge $i$. The vertex labels are
\begin{equation}
 \alpha_1=\epsilon_1,\qquad
 \alpha_v=\epsilon_{v-1}\oplus\epsilon_v\ (2\leq v\leq m),\qquad
 \alpha_{m+1}=\epsilon_m.
 \label{eq:vertex-labels}
\end{equation}
Thus the all-bottom problem has $2^{n-3}$ component blocks. The ordered
vertices are the left cap $\rho_{\alpha_1}(\phi_{h_1},\phi_{d_2},\phi_{d_1})$,
the interior forms
$\rho_{\alpha_v}(\phi_{h_v},\phi_{d_{v+1}},\phi_{h_{v-1}})$,
and the right cap
$\rho_{\alpha_{m+1}}(\phi_{d_n},\phi_{d_{n-1}},\phi_{h_m})$,
with descendants inserted as appropriate. Primary structure constants
are stripped off.

\section{The pillow map and the sphere conformal factor}
\label{sec:geometry}

\subsection{Periods, the normalized elliptic integral, and the coordinate}

Use the double cover and holomorphic differential
\begin{equation}
 y^2=x(x-z)(x-1),\qquad \omega=\frac{dx}{y}.
\end{equation}
With the cycle orientations chosen to give $\operatorname{Im}\tau>0$,
the periods and the human-note normalization of $K$ are
\begin{align}
 \oint_A\omega&=2\pi K(z),&
 \oint_B\omega&=2\pi iK(1-z),\notag\\
 K(z)&=\frac1\pi\int_0^1
 \frac{d\xi}{\sqrt{\xi(1-\xi)(1-z\xi)}}
 ={}_2F_1\!\left(\frac12,\frac12;1;z\right),
 \label{eq:K}\\
 \tau&=i\frac{K(1-z)}{K(z)},&\fq&=e^{\pi i\tau}.
 \label{eq:nome}
\end{align}
In particular, $K(0)=1$. Set
$\theta_j(\tau)=\vartheta_j(0\mid\tau)$. When a nome is used as the
argument, $\theta_3(\fq)=\sum_{k\in\mathbb Z}\fq^{k^2}$, so
$\theta_3(\fq)=\theta_3(\tau)$ and
$\theta_3(\fq^2)=\vartheta_3(0\mid2\tau)$. Then
\begin{equation}
 z=\frac{\theta_2(\tau)^4}{\theta_3(\tau)^4},\qquad
 K(z)=\theta_3(\tau)^2,\qquad
 w(x)=\frac1{\theta_3(\tau)^2}
       \int_0^x\frac{d\xi}{\sqrt{\xi(\xi-z)(\xi-1)}}.
 \label{eq:AJ}
\end{equation}
Thus $w\sim w+2\pi\sim w+2\pi\tau$. On the chosen sheet the branch
points map to $0,\pi,\pi+\pi\tau,\pi\tau$, respectively. The quotient
$w\sim-w$ is the pillow. Write $w_j=w(t_j)$ for the mobile insertions.

\subsection{Propagation factors and their branches}

For $n\geq5$, define the OPE-aligned cylinder parameters and effective
propagation factors by
\begin{align}
 p_1&=-e^{iw_1},&\varrho_1&=4p_1,\notag\\
 p_i&=e^{i(w_i-w_{i-1})}\quad(2\leq i\leq m-1),
 &\varrho_i&=p_i,\notag\\
 p_m&=-\fq e^{-iw_{m-1}},&\varrho_m&=4p_m.
 \label{eq:propagation}
\end{align}
They obey
\begin{equation}
 \prod_{i=1}^m p_i=\fq,\qquad
 \prod_{i=1}^m\varrho_i=16\fq.
 \label{eq:propagation-product}
\end{equation}
For $n=4$, set $p_1=\fq$ and $\varrho_1=16\fq$ instead.
We use $\varrho_i$ to distinguish propagation factors from the three-form
$\rho_\alpha$.

The two endpoint factors of $4$ encode the cap normalization, and each
interior segment has no such factor. Their product gives the same
$16\fq$ as the four-point OPE normalization. The endpoint minus signs
in \eqref{eq:propagation} are fixed by the aligned sheet: there
$w_j=\pi+i v_j$ and all $p_i$ are positive. Take positive square roots
there and continue them coherently. Since NS null levels can be
half-integral, these signs cannot be discarded or reselected at recursive
calls. They are geometric branch choices, not the edge-parity labels
$\epsilon_i$.

\subsection{Geometric prefactor and normalized pillow block}

The flat metrics are related by
\begin{equation}
 dw\,d\bar w=e^{2\sigma}dx\,d\bar x,\qquad
 e^{2\sigma}=\frac1{|\theta_3(\tau)|^4|x(x-z)(x-1)|}.
\end{equation}
A mobile bottom primary transforms as
$\phi_x(t)=(dw/dt)^d\phi_w(w(t))$. Its chiral Jacobian is therefore
\begin{equation}
 \left(\frac{dw}{dt}\right)^d
 =\theta_3(\tau)^{-2d}[t(t-z)(t-1)]^{-d/2}.
 \label{eq:Jacobian}
\end{equation}
On the ordered real sheet we absorb the fixed local phase into the
insertion state and write the positive radicand $t(1-t)(t-z)$ below.
Combining these Jacobians with the corner factors and Weyl anomaly gives
the scalar coordinate factor
\begin{equation}
 \begin{aligned}
 \Lambda_n^{(c)}={}&
 z^{c/24-d_1-d_2}
 (1-z)^{c/24-d_2-d_{n-1}}\\
 &\times\theta_3(\tau)^{c/2
       -4(d_1+d_2+d_{n-1}+d_n)-2\sum_{j=1}^{n-4}d_{j+2}}\\
 &\times\prod_{j=1}^{n-4}
       [t_j(1-t_j)(t_j-z)]^{-d_{j+2}/2}.
 \end{aligned}
 \label{eq:Lambda}
\end{equation}
This is the same bottom-primary coordinate factor as in the bosonic
construction; supersymmetry changes
the sewn modules, not the ordinary primary Jacobians.

Following the human note, $\cG$ denotes the pillow cap matrix element
after removing the primary propagation powers. Its sphere reconstruction is
\begin{equation}
 \boxed{\quad
 \cF^{\NS}_{c,\ve}(\vd,\vh;z,\boldsymbol t)
 =\Lambda_n^{(c)}
   \prod_{i=1}^m\varrho_i^{h_i-c/24}\,
   \cG_{n,\ve}(\vh;\boldsymbol p).
 \quad}
 \label{eq:sphere-G}
\end{equation}
At four points this is the parity-resolved version of the human note's
definition:
\begin{equation}
 \cG_{4,\epsilon}
 =(16\fq)^{-h+c/24}
 \langle\psi_{43}|\mathbb P_h\,
       \fq^{L_0-c/24}|\psi_{12}\rangle_{\epsilon}.
\end{equation}
For the even primary term the cap overlaps multiply to $16^{h-c/24}$,
as in the human note. The odd component has no level-zero term.

\section{The regular part: the NS pillow product}
\label{sec:seed}

Define
\begin{equation}
 P(\fq)=\prod_{j\geq1}(1-\fq^{2j}).
\end{equation}
In the human note's normalization, the known four-point all-bottom
large-weight limits are~\cite{HJS}
\begin{equation}
 \lim_{h\to\infty}\cG_{4,0}(h;\fq)
 =\frac{\theta_3(\fq^2)}{P(\fq)^{3/4}},\qquad
 \lim_{h\to\infty}\cG_{4,1}(h;\fq)=0.
 \label{eq:four-seed}
\end{equation}
The theta identity useful for checking the sphere normalization is
\begin{equation}
 z(1-z)\theta_3(\tau)^{12}=16\fq\,P(\fq)^{12}.
 \label{eq:theta-identity}
\end{equation}

Use the human note's common-weight variable,
\begin{equation}
 a_1=0,\qquad a_i=h_i-h_1\quad(2\leq i\leq m).
\end{equation}
The multipoint proposal is
\begin{equation}
 \boxed{\quad
 \lim_{\substack{h_1\to\infty\\a_i\ \mathrm{fixed}}}
 \cG_{n,\ve}
 =\delta_{\ve,\boldsymbol0}\,\cC_{\NS}(\fq),\qquad
 \cC_{\NS}(\fq)
 =\frac{\theta_3(\fq^2)}{P(\fq)^{3/4}}.
 \quad}
 \label{eq:seed}
\end{equation}
By the Jacobi triple product, this is explicitly
\begin{equation}
 \boxed{\quad
 \cC_{\NS}(\fq)
 =\prod_{j=1}^{\infty}
 \frac{(1-\fq^{4j})(1+\fq^{4j-2})^2}
      {(1-\fq^{2j})^{3/4}}.
 \quad}
 \label{eq:NS-product}
\end{equation}
Fractional powers are continued from $\fq=0$, where the product is one.
Equation~\eqref{eq:NS-product} replaces the bosonic pillow product
$P(\fq)^{-1/2}$ while keeping the geometric factor
$\Lambda_n^{(c)}$ unchanged. The dependence of the regular part on the
mobile insertion positions is only through
$\fq=\prod_i p_i$. The propagation powers in \eqref{eq:sphere-G}
retain the individual-position dependence.

The proposed large-weight reduction is that leading normalized even
interior vertices transport the same oscillator state, leaving the
four-point cap contribution. Parity-changing vertices have no leading
contribution in the present normalization. This motivates
\eqref{eq:seed}; a direct cap Ward-identity proof is not supplied here.

The pillow factor must not be identified with an unprojected torus trace.
The ordinary NS descendant character at torus nome $\fq^2$ is
$\prod_{j\geq1}(1+\fq^{2j-1})/(1-\fq^{2j})$. It has a term linear in
$\fq$, whereas the cap seed in \eqref{eq:four-seed} is
\begin{equation}
 \cC_{\NS}(\fq)=1+\frac{11}{4}\fq^2
                  +\frac{93}{32}\fq^4+O(\fq^6).
\end{equation}

\section{Singular part and the closed recursion}
\label{sec:recursion}

\subsection{Null-vector data in the human-note convention}

At fixed generic $c$, the NS Kac weights, levels, and null parities are
\begin{equation}
 h_{rs}(c)=\frac{Q^2-(rb+sb^{-1})^2}{8},\qquad
 \ell_{rs}=\frac{rs}{2},\qquad \pi_{rs}=rs\bmod2,
 \qquad r,s\geq1,\quad r+s\in2\mathbb Z.
 \label{eq:Kac}
\end{equation}
The range includes $(r,s)=(1,1)$, for which $h_{11}=0$ and
$\ell_{11}=1/2$. The $r\geq2$ restriction used to solve for a moving
$c$-pole does not apply to this fixed-$c$ $h$-recursion.

The inverse null-norm slope is the human note's
\begin{equation}
 \mathbf A^c_{rs}
 =\frac12
 \prod_{\substack{u=1-r,\ldots,r\\v=1-s,\ldots,s\\
                    u+v\in2\mathbb Z\\
                    (u,v)\ne(0,0),(r,s)}}
 \left[\frac{ub+vb^{-1}}{\sqrt2}\right]^{-1}.
 \label{eq:A}
\end{equation}
Here $b$ is fixed by $c$, not recalculated from a shifted internal weight.
For any weight $x$, set $\lambda_x^2=Q^2-8x$. Define
\begin{equation}
 \mathbf P^{\alpha}_{rs}(x,y;c)
 =\prod_{\substack{u=1-r,3-r,\ldots,r-1\\
                    v=1-s,3-s,\ldots,s-1\\
                    u+v-r-s\in4\mathbb Z+2(1-\alpha)}}
 \frac{\lambda_x-\lambda_y+ub+vb^{-1}}{2\sqrt2}
 \frac{\lambda_x+\lambda_y+ub+vb^{-1}}{2\sqrt2},
 \qquad \alpha\in\{0,1\}.
 \label{eq:P}
\end{equation}
These are precisely the $\mathbf P^0,\mathbf P^1$ of the human note.
For example $\mathbf P^0_{11}(x,y;c)=y-x$, $\mathbf P^1_{11}=1$,
and $\mathbf A^c_{11}=1/2$. Their argument order matters:
\begin{equation}
 \mathbf P^0_{rs}(x,y;c)=(-1)^{\lceil rs/2\rceil}
                         \mathbf P^0_{rs}(y,x;c),\qquad
 \mathbf P^1_{rs}(x,y;c)=(-1)^{\lfloor rs/2\rfloor}
                         \mathbf P^1_{rs}(y,x;c).
\end{equation}

For a pole on edge $k$, define the pole values and shifted child by
\begin{align}
 h_j^*&=h_j-h_k+h_{rs}=a_j-a_k+h_{rs},\notag\\
 h'_j&=h_j^*+\ell_{rs}\delta_{jk},&
 \epsilon'_j&=\epsilon_j\oplus\bigl(\pi_{rs}\delta_{jk}\bigr).
 \label{eq:child}
\end{align}
The two weight pairs adjacent to edge $k$ are
\begin{equation}
 (x_L,y_L)=
 \begin{cases}(d_1,d_2),&k=1,\\
              (h_{k-1}^*,d_{k+1}),&k>1,
 \end{cases}
 \qquad
 (x_R,y_R)=
 \begin{cases}(d_n,d_{n-1}),&k=m,\\
              (h_{k+1}^*,d_{k+2}),&k<m.
 \end{cases}
 \label{eq:adjacent-pairs}
\end{equation}
In the oriented bottom-component convention of
Section~\ref{sec:conventions}, the scalar residue kernel is
\begin{equation}
 \boxed{\quad
 \mathbf R^{(k,\ve)}_{rs}(\vh^*)
 =(-1)^{rs}\mathbf A^c_{rs}
   \mathbf P^{\alpha_k}_{rs}(x_L,y_L;c)
   \mathbf P^{\alpha_{k+1}}_{rs}(x_R,y_R;c).
 \quad}
 \label{eq:R}
\end{equation}
The sign is tied to the human note's fixed-parity three-form ordering,
not an extra universal sign to append to an independently normalized
formula. An odd null vector flips the edge label and both adjacent
vertex labels; an even null vector flips neither. This ordered kernel
is the one used in the numerical check described below.

\subsection{Regular part plus recursively factorized residues}

The common-weight pole is $h_1=h_{rs}-a_k$. At the pole,
\begin{equation}
 \Res_{h_1=h_{rs}-a_k}\cG_{n,\ve}
 =\varrho_k^{\ell_{rs}}
   \mathbf R^{(k,\ve)}_{rs}(\vh^*)\,
   \cG_{n,\ve'}(\vh').
 \label{eq:residue}
\end{equation}
The propagation factor follows by comparing the primary powers before
and after replacing the null state by its primary of weight
$h_{rs}+\ell_{rs}$:
\begin{equation}
 \frac{\prod_j\varrho_j^{h_j^*+\ell_{rs}\delta_{jk}-c/24}}
      {\prod_j\varrho_j^{h_j^*-c/24}}
 =\varrho_k^{\ell_{rs}}.
\end{equation}
The geometric factor is independent of internal weights and cancels.
The null-state matrix elements can therefore be evaluated in the plane
with the human-note fusion polynomials and transformed back.

At a fixed elliptic order, assume coefficientwise rationality in $h_1$
and generic, noncoincident simple poles. Subtracting their principal
parts leaves a polynomial in $h_1$. The finite large-$h_1$ limit in
\eqref{eq:seed} fixes that polynomial to the constant regular part.
Consequently the proposed recursion is
\begin{equation}
 \boxed{\begin{aligned}
 \cG_{n,\ve}(\vh)
 ={}&\underbrace{\delta_{\ve,\boldsymbol0}\cC_{\NS}(\fq)}
                         _{\text{regular part in }h_1}\\
 &+\sum_{k=1}^m\sum_{\substack{r,s\geq1\\r+s\in2\mathbb Z}}
 \frac{\varrho_k^{rs/2}\mathbf R^{(k,\ve)}_{rs}(\vh^*)}
      {h_k-h_{rs}(c)}\,
 \cG_{n,\ve'}(\vh').
 \end{aligned}}
 \label{eq:G-recursion}
\end{equation}
The sum is the singular part. There is no separate $o(1)$ remainder:
finite-weight corrections are supplied by the pole terms. For the
all-even component the regular part is the infinite product, and for
the other components it is zero. ``Regular'' refers to internal-weight
dependence, not to an expansion order in $\fq$.

\subsection{The unit-seed elliptic block and reconstruction}

As in the human note, reserve $H$ for the block with the regular product
removed:
\begin{equation}
 H_{n,\ve}:=\cC_{\NS}(\fq)^{-1}\cG_{n,\ve}.
 \label{eq:unit-H}
\end{equation}
It obeys
\begin{equation}
 \boxed{\quad
 H_{n,\ve}(\vh)=\delta_{\ve,\boldsymbol0}
 +\sum_{k=1}^m\sum_{\substack{r,s\geq1\\r+s\in2\mathbb Z}}
 \frac{\varrho_k^{rs/2}\mathbf R^{(k,\ve)}_{rs}(\vh^*)}
      {h_k-h_{rs}(c)}\,
 H_{n,\ve'}(\vh').
 \quad}
 \label{eq:H-recursion}
\end{equation}
All external weights, $c$, and coordinates remain fixed. The internal
differences for the next call are obtained from the child tuple
$\vh'$; they are not frozen to their parent values. Since
$\cC_{\NS}$ is weight independent, there is no additional ratio of
theta functions in the kernel. The final sphere block is
\begin{equation}
 \boxed{\quad
 \cF^{\NS}_{c,\ve}
 =\Lambda_n^{(c)}
   \prod_i\varrho_i^{h_i-c/24}\,
   \cC_{\NS}(\fq)\,H_{n,\ve}.
 \quad}
 \label{eq:reconstruction}
\end{equation}

\section{Four- and five-point specializations}
\label{sec:examples}

For four points, $\varrho_1=16\fq$ and
$\alpha_1=\alpha_2=\epsilon$. Thus
\begin{equation}
 \begin{aligned}
 H_{4,\epsilon}(h)={}&\delta_{\epsilon,0}\\
 &+\sum_{\substack{r,s\geq1\\r+s\in2\mathbb Z}}
 \frac{(16\fq)^{rs/2}(-1)^{rs}\mathbf A^c_{rs}
       \mathbf P^\epsilon_{rs}(d_1,d_2;c)
       \mathbf P^\epsilon_{rs}(d_4,d_3;c)}{h-h_{rs}}
 H_{4,\epsilon\oplus\pi_{rs}}(h_{rs}+rs/2).
 \end{aligned}
 \label{eq:four-recursion}
\end{equation}
The odd component is understood in the human-note phase convention.
As a useful sign check, the $(1,1)$ term
gives $H_{4,1}=-2\fq^{1/2}/h+\cdots$; in the plane this corresponds
to the leading coefficient $-1/(2h)$ multiplying
$z^{h-d_1-d_2+1/2}$ in the same convention.

For five points, the weights at $(0,z,t,1,\infty)$ are
$(d_1,d_2,d_3,d_4,d_5)$ and
\begin{equation}
 \varrho_1=-4e^{iw(t)},\qquad
 \varrho_2=-4\fq e^{-iw(t)},\qquad a=h_2-h_1.
\end{equation}
For each $(r,s)$ define
\begin{align}
 \mathbf R^{(1)}_{rs}
 &=(-1)^{rs}\mathbf A^c_{rs}
   \mathbf P^{\epsilon_1}_{rs}(d_1,d_2;c)
   \mathbf P^{\epsilon_1\oplus\epsilon_2}_{rs}(h_{rs}+a,d_3;c),
 \notag\\
 \mathbf R^{(2)}_{rs}
 &=(-1)^{rs}\mathbf A^c_{rs}
   \mathbf P^{\epsilon_1\oplus\epsilon_2}_{rs}(h_{rs}-a,d_3;c)
   \mathbf P^{\epsilon_2}_{rs}(d_5,d_4;c).
\end{align}
The complete two-edge recursion is
\begin{equation}
 \boxed{\begin{aligned}
 H_{5,\epsilon_1\epsilon_2}(h_1,h_2)
 ={}&\delta_{\epsilon_1,0}\delta_{\epsilon_2,0}\\
 &+\sum_{\substack{r,s\geq1\\r+s\in2\mathbb Z}}
 \frac{\varrho_1^{rs/2}\mathbf R^{(1)}_{rs}}{h_1-h_{rs}}
 H_{5,\epsilon_1\oplus\pi_{rs},\epsilon_2}
       (h_{rs}+rs/2,h_{rs}+a)\\
 &+\sum_{\substack{r,s\geq1\\r+s\in2\mathbb Z}}
 \frac{\varrho_2^{rs/2}\mathbf R^{(2)}_{rs}}{h_2-h_{rs}}
 H_{5,\epsilon_1,\epsilon_2\oplus\pi_{rs}}
       (h_{rs}-a,h_{rs}+rs/2).
 \end{aligned}}
 \label{eq:five-recursion}
\end{equation}
Its full sphere prefactor is \eqref{eq:Lambda} with $n=5$, and
$\cC_{\NS}$ is unchanged. No new function of $t$ is inserted into the
regular part.

\section{Implementation, checks, and status}
\label{sec:checks}

To evaluate a component block, compute $(\fq,w_j,p_i,\varrho_i)$ on a
coherent sheet, form the labels \eqref{eq:vertex-labels}, and apply
\eqref{eq:H-recursion} using \eqref{eq:A}--\eqref{eq:R}. Each recursive
term lowers the requested level on edge $k$ by $rs/2$, so a finite-order
coefficient calculation terminates. A twice-level integer index avoids
floating-point comparisons of half-integers. Finally multiply by
\eqref{eq:reconstruction}. Coincident Kac poles or coincident common-weight
pole locations require limits or a separate confluent treatment.

\subsection{Independent PBW audit: definition of the comparison}

The direct calculation uses only the NS algebra \eqref{eq:algebra}, the
primary Ward identities, and the three-form normalization
\eqref{eq:three-form-convention}. It does not call either $c$- or
$h$-recursion, or use Kac weights or fusion polynomials. At each level
$\ell_i$ it constructs the PBW Gram matrix $G_i$ and the ordered
three-point matrices. The plane coefficient is their contraction,
\begin{equation}
 B_{\boldsymbol\ell}
 =\beta_R^{\mathsf T}G_m^{-1}V_{m-1}G_{m-1}^{-1}\cdots
     V_1G_1^{-1}\beta_L.
 \label{eq:PBW-coefficient}
\end{equation}
Here $V_i$ has bra weight $h_{i+1}$, inserted weight $d_{i+2}$, and
ket weight $h_i$; $\beta_L,\beta_R$ are the cap vectors in the ordering
specified in Section~\ref{sec:conventions}. For one edge this reduces to
$\beta_R^{\mathsf T}G_1^{-1}\beta_L$.
All parity sectors are retained, with $\epsilon_i=2\ell_i\bmod2$.

Set $x_1=z/t_1$, $x_i=t_{i-1}/t_i$ for $2\leq i<m$, and
$x_m=t_{m-1}$; for four points use $x_1=z$. Then the plane expansion is
\begin{equation}
 \cF^{\NS}_{c,\ve}
 =z^{h_1-d_1-d_2}
   \prod_{j=1}^{m-1}t_j^{h_{j+1}-h_j-d_{j+2}}
   \sum_{\ell_i\in\mathbb Z_{\geq0}+\epsilon_i/2}
      B_{\boldsymbol\ell}\prod_i x_i^{\ell_i}.
 \label{eq:PBW-plane-series}
\end{equation}
We substitute the pillow map and divide by the entire prefactor in
\eqref{eq:reconstruction}; thus the quantity tested is precisely $H$,
not a block in a different normalization.

For reproducibility, the coordinate expansion is generated to the
requested order from products, rather than from a fixed low-order
polynomial. Write $U_j=\prod_{i\leq j}p_i$,
$V_j=\prod_{i>j}p_i$, so $U_jV_j=\fq$, and define
\begin{align}
 Z:=\frac{z}{16\fq}
 &=\prod_{k\geq1}\frac{(1+\fq^{2k})^8}{(1+\fq^{2k-1})^8},
 \label{eq:audit-Z}\\
 Y_j:=\frac{t_j}{4V_j}
 &=(1+U_j)^2\prod_{k\geq1}
 \frac{(1+\fq^{2k})^4}{(1+\fq^{2k-1})^4}
 \frac{(1+U_j\fq^{2k})^2(1+V_j\fq^{2k-1})^2}
      {(1+U_j\fq^{2k-1})^2(1+V_j\fq^{2k-2})^2}.
 \label{eq:audit-Y}
\end{align}
The coordinate units are $x_1/(4p_1)=Z/Y_1$,
$x_i/p_i=Y_{i-1}/Y_i$ in the interior, and $x_m/(4p_m)=Y_{m-1}$.
Their constant terms are one. The multiplier applied to the composed
plane descendant series, after the primary powers cancel, is
\begin{equation}
 \begin{split}
 &Z^{h_1-c/24}\prod_{j=1}^{m-1}Y_j^{h_{j+1}-h_j}
 \prod_{j=1}^{m-1}\bigl[(1-t_j)(1-z/t_j)\bigr]^{d_{j+2}/2}\\
 &\quad\times(1-z)^{-c/24+d_2+d_{n-1}}
 \theta_3(\fq)^{-c/2+4(d_1+d_2+d_{n-1}+d_n)
                        +2\sum_jd_{j+2}}
 \cC_{\NS}(\fq)^{-1}.
 \end{split}
 \label{eq:audit-multiplier}
\end{equation}
Consequently the check tests the coordinate map, all conformal factors,
the fermionic signs, the regular product, and the residue recursion
together. No effective shift of $c$ is used in this implementation.

``Through degree $N$'' means every coefficient of
$\prod_i p_i^{\ell_i}$ with $\ell_i\geq0$ half-integral and
$\sum_i\ell_i\leq N$. It is not a rectangular cutoff on each edge.
The number of tested coefficients is $\binom{2N+m}{m}$, including the
constant and the vanishing coefficients. At $N=10$ the largest individual
Gram matrix has dimension $161$.

As an independent algebra check, the new PBW/Ward implementation was
compared with the existing exact descendant-three-form implementation
through total level three: all $57$ Gram entries and $110$ Ward entries
agree identically for independent symbols $c,h,d,k$.

\subsection{Analytic PBW comparison through total degree three}

All parameters are kept symbolic, with coefficients in
$\mathbb Q(b,d_1,\ldots,d_n,h_1,\ldots,h_m)$. The PBW calculation first
uses an independent $c$, and only at the comparison stage substitutes
$c=3/2+3(b+b^{-1})^2$. Thus an exact comparison here does not mean
evaluation at a special rational parameter set.

There are two equivalent symbolic tests. One expands both sides of
\eqref{eq:H-recursion} and reduces every coefficient difference to zero.
The other gives a smaller rational-identity certificate. Write
$C_{\boldsymbol\ell}=[\prod_i p_i^{\ell_i}]H$ and
$\kappa_i=\varrho_i/p_i$. For every coefficient computed by PBW, check:
\begin{enumerate}
\item Its reduced denominator, viewed over $\mathbb Q(b)$, divides
\begin{equation}
 \prod_{k=1}^m\prod_{\substack{r,s\geq1,\ r+s\in2\mathbb Z\\rs\leq2\ell_k}}
      (h_k-h_{rs}).
 \label{eq:audit-pole-catalog}
\end{equation}
Thus no additional poles or higher-order poles are left unchecked.
\item Every residue equals the proposed kernel times the lower-degree
PBW coefficient:
\begin{equation}
 \Res_{h_k=h_{rs}} C_{\boldsymbol\ell}
 =\kappa_k^{rs/2}
    \left.\mathbf R^{(k,\ve)}_{rs}\right|_{h_k=h_{rs}}
    \left.C_{\boldsymbol\ell-(rs/2)\boldsymbol e_k}
                    \right|_{h_k=h_{rs}+rs/2}.
 \label{eq:audit-residue-identity}
\end{equation}
In this local test the other $h_j$ are independent symbolic variables
held fixed. Substitution of their pole values $h_j^*$ then gives exactly
the common-weight residue in \eqref{eq:H-recursion}.
\item The directly computed regular part is
\begin{equation}
 \lim_{L\to\infty}C_{\boldsymbol\ell}(h_1+L,\ldots,h_m+L)
 =\delta_{\boldsymbol\ell,\boldsymbol0}.
 \label{eq:audit-regular-part}
\end{equation}
This is tested from numerator and denominator degrees in $L$, including
their leading coefficients; the seed is not supplied to the PBW side.
\end{enumerate}
With the differences $h_j-h_1$ fixed, the difference from the recursive
answer is a rational function with no poles and zero regular part.
It therefore vanishes identically. Induction in total degree makes these
three checks an exact coefficient comparison with the expanded recursion,
not merely a check of its pole locations.

The completed generic-symbolic audit gives
\begin{center}
\begin{tabular}{@{}ccll@{}}
\toprule
Points & Coefficients & Exact method & Result\\
\midrule
4 & 7 & Expanded coefficient differences & All zero\\
5 & 28 & Expanded differences and certificate & All zero\\
6 & 84 & Pole and regular-part certificate & All zero\\
\bottomrule
\end{tabular}
\end{center}
The five-point certificate checks all $106$ candidate residue identities,
and the six-point certificate all $342$, including vanishing residues.
Every coefficient also passes the denominator-divisibility and
large-common-weight regular-part tests. There is no numerical tolerance
or specialization of $b,d_i,h_i$ in these results. No change to the
proposed recursion or its conformal factor is required by this audit.

For example, both independent calculations give
\begin{align}
 H_{5,00}={}&1
 -\frac{2(d_1-d_2)(d_3+h_1-h_2)}{h_1}\,p_1
 +\frac{2(d_4-d_5)(d_3-h_1+h_2)}{h_2}\,p_2
 +O_{\rm tot}(2),\notag\\
 H_{5,10}={}&-\frac{\sqrt{p_1}}{h_1}+O_{\rm tot}(3/2),\qquad
 H_{5,01}=-\frac{\sqrt{p_2}}{h_2}+O_{\rm tot}(3/2),\notag\\
 H_{5,11}={}&\frac{d_3-h_1-h_2}{h_1h_2}\sqrt{p_1p_2}
                  +O_{\rm tot}(2).
 \label{eq:audit-leading-coefficients}
\end{align}
Here $O_{\rm tot}(a)$ contains terms of total degree at least $a$.
These signs refer throughout to the human-note three-form convention.
The six-point leading terms additionally distinguish an interior edge
from a cap edge:
\begin{align}
 H_{6,010}&=-\frac{\sqrt{p_2}}{2h_2}+O_{\rm tot}(3/2),\notag\\
 H_{6,101}&=\frac{\sqrt{p_1p_3}}{h_1h_3}+O_{\rm tot}(2),\notag\\
 H_{6,111}&=-\frac{(d_3-h_1-h_2)(d_4-h_2-h_3)}{2h_1h_2h_3}
              \sqrt{p_1p_2p_3}+O_{\rm tot}(5/2).
 \label{eq:audit-six-leading-coefficients}
\end{align}
In particular, the interior edge uses $\varrho_2=p_2$, not $4p_2$.

\subsection{Numerical PBW comparison through total degree ten}
\label{sec:PBW-numerical}

Four generic parameter sets were tested for both five and six points.
The six-point fixtures are listed below; the five-point fixture deletes
the fourth external entry and the third internal entry. In this table
$c=3/2+3(b+b^{-1})^2$, and the labels A--D refer to the PBW audit, not
the earlier degree-four fixtures below. The new degree-ten
$c$-recursion audit uses these same A--D fixtures.
\begin{center}
\small
\begin{tabular}{@{}ccll@{}}
\toprule
Case & $b$ & $(d_1,d_2,d_3,d_4,d_5,d_6)$ & $(h_1,h_2,h_3)$\\
\midrule
A & $1.27$ & $(.31,.42,.53,.37,.47,.28)$ & $(.73,1.10,1.37)$\\
B & $1.43$ & $(.22,.61,.39,.58,.74,.45)$ & $(1.13,.85,1.69)$\\
C & $.83$ & $(.17,.83,1.21,.46,.34,.67)$ & $(.19,1.47,.82)$\\
D & $1.61$ & $(1.12,.26,.79,1.33,.58,.91)$ & $(2.31,.64,1.09)$\\
\bottomrule
\end{tabular}
\end{center}
For each set, every coefficient through total degree ten was compared:
$231$ coefficients in all four parity sectors for five points, and
$1771$ coefficients in all eight sectors for six points. Using 80-digit
arithmetic, the maximum scaled errors are
\begin{center}
\begin{tabular}{@{}ccc@{}}
\toprule
PBW case & Five points & Six points\\
\midrule
A & $4.84\times10^{-74}$ & $1.80\times10^{-69}$\\
B & $9.34\times10^{-72}$ & $7.81\times10^{-69}$\\
C & $1.94\times10^{-71}$ & $7.57\times10^{-69}$\\
D & $3.37\times10^{-70}$ & $2.42\times10^{-67}$\\
\bottomrule
\end{tabular}
\end{center}
Let $u_{\boldsymbol\ell}$ and $v_{\boldsymbol\ell}$ denote the coefficient
of the same $H$ computed by PBW and recursion, respectively. The error is
\begin{equation}
 \frac{|u_{\boldsymbol\ell}-v_{\boldsymbol\ell}|}
 {\max\{1,|u_{\boldsymbol\ell}|,|v_{\boldsymbol\ell}|\}}.
 \label{eq:audit-error}
\end{equation}
Every coefficient passes
the $10^{-50}$ tolerance. Individual edge levels can reach ten; the
comparison does not restrict each edge to the evenly shared cutoff.
The general numerical audit includes the four-point control at case A:
all $21$ coefficients through degree ten agree with maximum scaled
error $7.01\times10^{-70}$.

Repeating six-point case C at 110-digit precision gives maximum scaled
error $1.48\times10^{-97}$ for the same $1771$ coefficients. Comparing
the 80- and 110-digit runs directly, the PBW coefficients change by at
most $3.04\times10^{-74}$, and the recursion coefficients by at most
$7.57\times10^{-69}$ in the same scaled norm. This precision improvement
is consistent with finite-precision cancellation rather than a residual
disagreement in the formula.

These tests include the regular product, whose first nonconstant term
occurs at total $p_i$-degree $2m$: four for five points and six for six
points. Thus the degree-ten audit is stronger than a low-order comparison
that never sees the proposed NS regular product.

\subsection{Reproducing the PBW audit}

The implementation lives in
\nolinkurl{Code/sphere_five_kummer_h_recursion/}.
The independent sewing engine is \nolinkurl{ns_pillow_direct_pbw.py};
the coordinate products and proposed recursion are in
\nolinkurl{ns_pillow_elliptic_audit.py}. From that directory run
\begin{quote}\small\ttfamily
python3 check\_ns\_pillow\_direct\_pbw.py\\
\hspace*{1em}-{}-mode symbolic -{}-points 5 -{}-degree 3\\[0.4em]
python3 check\_ns\_pillow\_direct\_pbw.py\\
\hspace*{1em}-{}-mode symbolic -{}-symbolic-method residues -{}-points 6 -{}-degree 3\\[0.4em]
python3 check\_ns\_pillow\_direct\_pbw.py\\
\hspace*{1em}-{}-mode numeric -{}-points 6 -{}-degree 10 -{}-case A -{}-dps 80
\end{quote}
Each indented line continues the command above it. Repeat the numerical
command for both five and six points and cases A--D. Use
\texttt{-{}-case C -{}-dps 110} for the precision repeat and
\texttt{-{}-mode oracle} for the independent low-level Gram/Ward audit.
Appending \texttt{-{}-output result.json} saves the coefficient ledger.
The saved audit files are under \texttt{Data Set/h-Recursion/}, with
names beginning \texttt{ns\_pillow\_pbw\_}. The numerical files contain
both coefficient values, precision, error by degree and parity, and the
minimum visited recursion denominator. The symbolic files retain an
exact-zero status for each coefficient and, for the certificate method,
the individual pole and regular-part checks.

The five quick regression tests also check signs, truncation consistency,
and rejection of deliberately incorrect residues, regular parts, and
extra poles. The numerical PBW fixtures have positive-definite Gram
matrices and use diagonally scaled Cholesky solves. This audit does not
claim to implement a general nonunitary or complex-weight Gram solver.

\subsection{Independent literature \texorpdfstring{$c$}{c}-recursion: conventions and seed}
\label{sec:c-literature}

The independent comparison uses the sphere linear-channel recursion of
Belavin--Geiko~\cite{BG}, equations (3.12)--(3.18), with the global
matrix elements of their Appendix~A. The four-point control uses their
equations (3.8)--(3.10). There is an indexing caveat in arXiv v1:
the interior line of (3.18) prints the adjacent external labels as
$d_{i+2},d_{i+3}$. With the $h_i$ ordering in their Figure~3 and
(3.15), the adjacent insertions are instead $d_{i+1},d_{i+2}$,
with the corresponding component labels. We use this graph-consistent
adjacency, not the literal shifted indices of that line. For six
points, for example, the middle edge joins insertions $d_3,d_4$.
This indexing correction is separate from a normalization change.

Their central charge, denoted $c_{\rm BG}$ here
only to make the translation explicit, satisfies
\begin{equation}
 c=\frac32c_{\rm BG},\qquad
 c_{rs}(h)=\frac32c^{\rm BG}_{rs}(h),\qquad
 J_{rs}(h):=-\frac{dc_{rs}(h)}{dh}
           =\frac32J^{\rm BG}_{rs}(h).
 \label{eq:c-charge-conversion}
\end{equation}
All subsequent formulas use only the human-note $c$. Converting a pole
but not its Jacobian would change the residue by a factor of $3/2$.
The three-form reflection phases and odd-null transport are likewise
translated into \eqref{eq:three-form-convention}; no second normalization
of $H$ is introduced.

The paper's parameters $q_i$ in its equation (3.16) are plane plumbing
ratios, not $\fq$ or the pillow $p_i$. After reversing its external and
internal ordering to \eqref{eq:positions}, they are the reverse-ordered
list of our $x_i$ in \eqref{eq:PBW-plane-series}. For example, at six points
their list becomes $(t_2,t_1/t_2,z/t_1)=(x_3,x_2,x_1)$.
Their normalized plane series is therefore compared with $H$ only after
restoring the primary powers and dividing by the entire prefactor in
\eqref{eq:reconstruction}, as in the PBW pullback above.

For clarity, the implemented coefficient recursion in our ordering is
\begin{equation}
 \begin{aligned}
 B_{\boldsymbol\ell}^{\ve}(c,\vh)
 ={}&B_{\boldsymbol\ell}^{\mathrm{glob},\ve}(\vh)\\
 &+\sum_{k=1}^{m}
 \sum_{\substack{r\geq2,\ s\geq1\\r+s\in2\mathbb Z\\rs\leq2\ell_k}}
 \frac{J_{rs}(h_k)\,\mathcal R_{rs}^{(k,\ve)}(\vh;c_*)}
      {c-c_*}\,
 B_{\boldsymbol\ell-(rs/2)\boldsymbol e_k}^{\ve'}
       \left(c_*,\vh+\frac{rs}{2}\boldsymbol e_k\right),
 \qquad c_*=c_{rs}(h_k).
 \end{aligned}
 \label{eq:literature-c-recursion}
\end{equation}
Here $\boldsymbol e_k$ is the $k$th unit vector,
$\epsilon'_j=\epsilon_j\oplus(\pi_{rs}\delta_{jk})$, and
$B_{\boldsymbol\ell}^{\ve}$ denotes the coefficient of the plane
descendant series with its leading powers removed. With the adjacent
pairs of \eqref{eq:adjacent-pairs}, now using the unshifted neighboring
$h_j$ rather than $h_j^*$, the residue is
\begin{equation}
 \mathcal R_{rs}^{(k,\ve)}(\vh;c_*)
 =(-1)^{rs}\mathbf A^{c_*}_{rs}
  \mathbf P^{\alpha_k}_{rs}(x_L,y_L;c_*)
  \mathbf P^{\alpha_{k+1}}_{rs}(x_R,y_R;c_*).
 \label{eq:c-residue}
\end{equation}
Every null-vector factor is evaluated at the moving central-charge pole.
One convenient branch, appropriate to the generic fixtures below, is
\begin{align}
 X_{rs}(h):=b_{rs}(h)^2
 &=-\frac{4h+rs-1+\sqrt{16h^2+8(rs-1)h+(r-s)^2}}{r^2-1},\notag\\
 c_{rs}(h)&=\frac{15}{2}+3\left(X_{rs}(h)+X_{rs}(h)^{-1}\right).
 \label{eq:c-pole}
\end{align}
The range $r\geq2$ is essential for this moving-$c$ parametrization.
Only $h_k$ shifts in its child, and $c$ changes to $c_*$. In contrast,
\eqref{eq:child} shifts the entire internal-weight tuple at fixed $c$.
Thus this is not a rearrangement of the implemented elliptic recursion.

The large-$c$ seed is a genuine global $\mathfrak{osp}(1|2)$ network.
Writing $\ell_i=N_i+\epsilon_i/2$ and
$v_i=L_{-1}^{N_i}G_{-1/2}^{\epsilon_i}\phi_{h_i}$, it is
\begin{equation}
 B_{\boldsymbol\ell}^{\mathrm{glob},\ve}
 =\frac{\rho_{\alpha_1}(v_1,\phi_{d_2},\phi_{d_1})
     \displaystyle\prod_{i=1}^{m-1}
       \rho_{\alpha_{i+1}}(v_{i+1},\phi_{d_{i+2}},v_i)
     \rho_{\alpha_{m+1}}(\phi_{d_n},\phi_{d_{n-1}},v_m)}
    {\displaystyle\prod_{i=1}^{m}N_i!\,(2h_i)_{N_i+\epsilon_i}}.
 \label{eq:c-global-seed}
\end{equation}
The matrix elements are evaluated algebraically in the global algebra,
with precisely the human-note ordering of $\rho_\alpha$.
This seed is neither the constant $1$ nor $\cC_{\NS}$.
In particular, the $c$-recursive plane calculation does not assume the
proposed large-internal-weight pillow seed.

The independent code is
\nolinkurl{Code/c_Recursion/ns_multipoint_c_recursion.py}, using
\nolinkurl{ns_recursion_recipe.py} for moving poles and scalar residues
and \nolinkurl{ns_global_osp_block.py} for the global seed. None calls
the elliptic recursion or a PBW solver. The new comparison driver only
adds memoization and selects the vertex labels for each edge-parity
tuple. The exact pillow map and geometric multiplier are shared with
the preceding PBW audit; the two plane-block constructions are not.

\subsection{Numerical literature comparison through total degree ten}
\label{sec:c-numerical}

The new test recomputes both the plane $c$-recursion and the elliptic
$h$-recursion at 80-digit precision, using all four fixtures in
Section~\ref{sec:PBW-numerical}. It compares their coefficients of the
same unit-seed $H$, with the scaled error \eqref{eq:audit-error}.
The cutoff is $\sum_i\ell_i\leq10$, including every half-integer edge
level: $231$ coefficients per five-point fixture in four parity sectors,
and $1771$ per six-point fixture in eight sectors. The maximum errors are
\begin{center}
\begin{tabular}{@{}ccc@{}}
\toprule
Case & Five points: $c$ versus $h$ & Six points: $c$ versus $h$\\
\midrule
A & $4.14\times10^{-74}$ & $1.80\times10^{-69}$\\
B & $9.34\times10^{-72}$ & $7.81\times10^{-69}$\\
C & $1.94\times10^{-71}$ & $7.57\times10^{-69}$\\
D & $3.37\times10^{-70}$ & $2.42\times10^{-67}$\\
\bottomrule
\end{tabular}
\end{center}
Every coefficient passes the $10^{-50}$ tolerance; no modification of
the proposed $h$-recursion or its prefactor is required. The four-point
control at case A also passes all $21$ coefficients, with maximum error
$1.41\times10^{-69}$.

As a third check, the newly computed $c$-recursive $H$ is compared with
the saved independent PBW coefficient ledgers. The maximum errors across
A--D are $1.47\times10^{-73}$ for five points and
$3.58\times10^{-74}$ for six points. These saved PBW values are used
only as comparison targets, never to construct either recursion.
Both recursions are computed afresh at each precision.

Repeating six-point case C at 110 digits gives maximum error
$1.48\times10^{-97}$ between the two recursions, and
$2.34\times10^{-104}$ between $c$-recursion and the independent PBW
ledger. Across the 80- and 110-digit runs, the transformed $c$-recursive
coefficients change by at most $1.28\times10^{-74}$ and the
$h$-recursive coefficients by $7.57\times10^{-69}$ in the same scaled
norm. The improvement is consistent with finite-precision cancellation,
not an order-ten discrepancy in the proposal.

The comparison includes the nonconstant terms of the NS product:
$\fq^2=(\prod_i p_i)^2$ first appears at total degree four for five
points and six for six points. Since the independent $c$-recursive seed
is \eqref{eq:c-global-seed}, the agreement checks the proposed pillow
regular product, rather than assuming it on both sides.

From \nolinkurl{Code/sphere_five_kummer_h_recursion/}, run
\begin{quote}\small\ttfamily
python3 check\_ns\_pillow\_c\_recursion.py\\
\hspace*{1em}-{}-points 6 -{}-case A -{}-degree 10 -{}-dps 80\\
\hspace*{1em}-{}-output result.json
\end{quote}
Repeat with five points and cases A--D; case C at 110 digits provides
the precision repeat. The optional \texttt{-{}-pbw-ledger} argument
takes the corresponding saved PBW JSON path and enables the third
comparison. It is not needed for $c$ versus $h$. The new files in
\texttt{Data Set/h-Recursion/} begin with
\texttt{ns\_pillow\_c\_comparison\_}. They retain every plane
$c$-recursive coefficient, its transformed $H$, the independently
computed $h$-recursive $H$, per-degree and per-parity errors, precision,
and source hashes. The summary JSON collects all runs and the precision
comparison. These are our numerical evaluations of the published
recursion, not numerical data supplied by the paper's authors.

Five additional regression tests check that memoization leaves the
unmodified literature implementation unchanged, exercise all six-point
sectors at degree three, verify the nonunit global seed and odd phase,
reject invalid level tuples, and detect an intentionally omitted
Jacobian conversion. The 13 existing multipoint $c$-recursion regression
tests also pass. The numerical test assumes generic nonconfluent moving
$c$-poles; it does not prescribe a limiting procedure at collisions.

\subsection{Earlier comparison with the NS \texorpdfstring{$c$}{c}-recursion}

The existing reproducible five-point test compares the elliptic proposal
against the independent multipoint NS $c$-recursion~\cite{BG}, pulled
back by the exact pillow coordinate series. It checks all four parity
channels through total elliptic degree four: 45 coefficients per parameter
set, including half-integer edge powers. This is not a rectangular
level-four check and not a direct PBW computation. The generic cases are
\begin{center}
\begin{tabular}{@{}ccll@{}}
\toprule
Case & $c$ & $(d_1,d_2,d_3,d_4,d_5)$ & $(h_1,h_2)$\\
\midrule
A & $14.7$ & $(0.31,0.42,0.53,0.47,0.28)$ & $(0.73,1.10)$\\
B & $21.3$ & $(0.22,0.61,0.39,0.74,0.45)$ & $(1.13,0.85)$\\
\bottomrule
\end{tabular}
\end{center}
At 75-digit working precision, with 60-digit comparison data, the maximum
relative errors were
\begin{center}
\begin{tabular}{@{}cc@{}}
\toprule
Case & Maximum relative error for $H$\\
\midrule
A & $2.7\times10^{-60}$\\
B & $1.9\times10^{-60}$\\
\bottomrule
\end{tabular}
\end{center}
The error denominator is $\max(1,|\text{left}|,|\text{right}|)$.
A separate large-weight check used $(h_1,h_1+0.37)$ in case A and
$(h_1,h_1-0.28)$ in case B, with $h_1=50,200,800$. The coefficientwise
deviations from the proposed seed decreased approximately as $h_1^{-1}$.
For example, the all-even $p_1^2p_2^2$ coefficient of
$H$ in case A was
$-0.0971729251,-0.0253007025,-0.0063915533$, approaching $0$.

The product identity in \eqref{eq:NS-product} was separately checked
symbolically through $\fq^{10}$. The prefactor conversion at
$\fq=0.01,0.1,0.3$ agreed to better than $2.3\times10^{-75}$.
Those separate tests by themselves check normalization; the actual
degree-ten block-coefficient comparisons are in
Sections~\ref{sec:PBW-numerical} and~\ref{sec:c-numerical}.
From \nolinkurl{Code/sphere_five_kummer_h_recursion/}, rerun the checks with:
\begin{quote}\small\ttfamily
python3 check\_ns\_sphere\_five\_elliptic\_h\_recursion.py
\end{quote}

The general-$n$ seed remains the proposal in \eqref{eq:seed}. The direct
PBW and independent $c$-recursion checks support the four-, five-, and six-point formulas at the
stated finite orders; they do not prove the seed for arbitrary $n$ or
control the remainder at a fixed nonzero coordinate. The bosonic
multipoint elliptic construction is prior work of Artemev and
Khromov~\cite{AK}, in particular their equations (3.16) and (5.14)--(5.18).
The four-point NS seed is prior work of Hadasz, Jask\'olski, and
Suchanek~\cite{HJS}, Section~6. We do not claim either as new, nor make a
novelty claim for this multipoint NS proposal without a separate
literature audit.

\begin{samepage}
\begin{thebibliography}{9}
\bibitem{AK}
A.~Artemev and D.~Khromov,
\emph{WKB-asymptotics for multipoint Virasoro conformal blocks and
applications}, \href{https://arxiv.org/abs/2603.08194v2}{arXiv:2603.08194v2}.

\bibitem{HJS}
L.~Hadasz, Z.~Jask\'olski, and P.~Suchanek,
\emph{Elliptic recurrence representation of the $N=1$ Neveu--Schwarz
blocks}, \href{https://arxiv.org/abs/0711.1619}{arXiv:0711.1619}.

\bibitem{BG}
V.~A.~Belavin and R.~V.~Geiko,
\emph{$c$-Recursion for multi-point superconformal blocks. NS sector},
\href{https://arxiv.org/abs/1806.09563}{arXiv:1806.09563}.
\end{thebibliography}
\end{samepage}

\end{document}
